\slcol{श्रीभगवानुवाच —\\
\Index{भूय एव महाबाहो} शृणु मे परमं वचः । \\
यत्तेऽहं प्रीयमाणाय वक्ष्यामि हितकाम्यया ॥ १०.१ ॥ }
\translate{śrībhagavānuvāca —\\
bhūya ēva mahābāhō śr̥ṇu mē paramaṁ vacaḥ | \\
yattē:'haṁ prīyamāṇāya vakṣyāmi hitakāmyayā || 10.1 || }
\cquote{Sri Bhagavan said: O Mahābāho (mighty-armed one, Arjuna), listen once again to My Supreme word. Because you are dear to me, I shall speak these words for your benefit and well-being.}
\slcol{\Index{न मे विदुः सुरगणाः} प्रभवं न महर्षयः । \\
अहमादिर्हि देवानां महर्षीणां च सर्वशः ॥ १०.२ ॥ }
\translate{na mē viduḥ suragaṇāḥ prabhavaṁ na maharṣayaḥ | \\
ahamādirhi dēvānāṁ maharṣīṇāṁ ca sarvaśaḥ || 10.2 || }
\cquote{Neither the hosts of gods nor the great sages know My origin, for in every way, I am the source of the gods and the great sages.}

\newpage
\begin{mananam}{\mananamfont {Mananam Śloka - 2, 3}}
\normalfont\mananamtext What is my concept of God, or the force behind creation? What is the nature of the entity that sustains all creation and enables our internal systems to function in harmony? Could it also have a beginning like all of us? If so, would it also be subject to an end?
What is the status of other subsidiary gods and sages in relation to this Supreme God? How might I benefit from learning and gaining clarity on the nature of the Supreme God?
\end{mananam}
\sreflect
\begin{inspiration}{\mananamfont {Inspiration}}
\normalfont\mananamtext The One Universal Force, the source of all other forces, from which everything arises and into which all dissolves, is the ultimate reality worthy of our highest pursuit. How can we truly know or understand this entity if we remain confined to the human intellect’s analytical and deductive processes? This entity, beyond cause and effect, without birth or death, can only be known and realised through the guidance of the scriptures and qualified teachers. By embracing the right sources of knowledge, we remain connected with God as the source of this world, rather than getting lost in God’s projection.
\end{inspiration}
\newpage

\slcol{\Index{यो मामजमनादिं} च वेत्ति लोकमहेश्वरम् । \\
असंमूढः स मर्त्येषु सर्वपापैः प्रमुच्यते ॥ १०.३ ॥ }
\translate{yō māmajamanādiṁ ca vētti lōkamahēśvaram | \\
asaṁmūḍhaḥ sa martyēṣu sarvapāpaiḥ pramucyatē || 10.3 || }
\cquote{He who knows Me as unborn, without beginning, and the Supreme Lord of all worlds—such a one, who is undeluded among mortals, is liberated from all sins.}
\slcol{\Index{बुद्धिर्ज्ञानमसंमोहः} क्षमा सत्यं दमः शमः । \\
सुखं दुःखं भवोऽभावो भयं चाभयमेव च ॥ १०.४ ॥ }
\translate{buddhirjñānamasaṁmōhaḥ kṣamā \\
\hspace*{0.5cm}satyaṁ damaḥ śamaḥ | \\
sukhaṁ duḥkhaṁ bhavō:'bhāvō \\
\hspace*{0.5cm}bhayaṁ cābhayamēva ca || 10.4 || }
\cquote{Intelligence, knowledge, freedom from delusion, forgiveness, truthfulness, control of the senses, control of the mind, pleasure and pain, birth and death, fear and fearlessness.}
\slcol{\Index{अहिंसा समता तुष्टिस्तपो} दानं यशोऽयशः । \\
भवन्ति भावा भूतानां मत्त एव पृथग्विधाः ॥ १०.५ ॥ }
\translate{ahiṁsā samatā tuṣṭistapō dānaṁ yaśō:'yaśaḥ | \\
bhavanti bhāvā bhūtānāṁ matta ēva pr̥thagvidhāḥ || 10.5 || }
\cquote{Nonviolence, equanimity, contentment, austerity, charity, fame and infamy—all these diverse qualities of living beings arise from Me alone.}
\slcol{\Index{महर्षयः सप्त पूर्वे} चत्वारो मनवस्तथा । \\
मद्भावा मानसा जाता येषां लोक इमाः प्रजाः ॥ १०.६ ॥ }
\translate{maharṣayaḥ sapta pūrvē catvārō manavastathā | \\
madbhāvā mānasā jātā yēṣāṁ lōka imāḥ prajāḥ || 10.6 || }
\cquote{The seven great sages, the four other great sages before them, the ancient \emph{manus} (progenitors of mankind), who are endowed with my essence, were born from My mind. From them, all the creatures of this world have originated.}

\newpage
\begin{mananam}{\mananamfont {Mananam Śloka - 4, 5}}
\normalfont\mananamtext What are the qualities of a wise person? Does being spiritually inclined help us to develop these qualities? What else shapes someone who embodies these positive traits? 
What is the role of daily life practice and training in cultivating a loving, kind, self-controlled, and equanimous nature? Is such a disposition sustained by maintaining the right attitude? Can these qualities be attained by someone who lacks spiritual wisdom? If not, why not?
\end{mananam}
\sreflect
\begin{inspiration}{\mananamfont {Inspiration}}
\normalfont\mananamtext All positive qualities find their perfection in the Lord, whom all spiritually inclined people seek and revere. To truly understand the pure essence of any of these qualities is to know the Lord. Conversely, ignorance of the Lord prevents the realisation of pure love, calmness, happiness, and other virtues. Knowing the Lord does not mean experiencing a supernatural vision of the Lord in a specific form, but recognising the essence of Being itself—the source of all creation. Gaining wisdom about the Lord and His qualities is the most effective way to cultivate, strengthen, and sustain all positive qualities in life.
\end{inspiration}
\newpage

\slcol{\Index{एतां विभूतिं योगं} च मम यो वेत्ति तत्त्वतः । \\
सोऽविकम्पेन योगेन युज्यते नात्र संशयः ॥ १०.७ ॥ }
\translate{ētāṁ vibhūtiṁ yōgaṁ ca mama yō vētti tattvataḥ | \\
sō:'vikampēna yōgēna yujyatē nātra saṁśayaḥ || 10.7 || }
\cquote{He who knows My divine glories (\emph{vibhūti}) and My \emph{yoga} power in their true essence, achieve unwavering union through \emph{yoga}—there is no doubt about this.}
\slcol{\Index{अहं सर्वस्य प्रभवो} मत्तः सर्वं प्रवर्तते । \\
इति मत्वा भजन्ते मां बुधा भावसमन्विताः ॥ १०.८ ॥ }
\translate{ahaṁ sarvasya prabhavō mattaḥ sarvaṁ pravartatē | \\
iti matvā bhajantē māṁ budhā bhāvasamanvitāḥ || 10.8 || }
\cquote{I am the source of all creation; from Me everything emanates; knowing thus, the wise endowed with devotional feeling (\emph{bhāva}) worship Me.}
\slcol{\Index{मच्चित्ता मद्गतप्राणा} बोधयन्तः परस्परम् । \\
कथयन्तश्च मां नित्यं तुष्यन्ति च रमन्ति च ॥ १०.९ ॥ }
\translate{maccittā madgataprāṇā bōdhayantaḥ parasparam | \\
kathayantaśca māṁ nityaṁ tuṣyanti ca ramanti ca || 10.9 || }
\cquote{With minds absorbed on Me and senses devoted to Me, mutually sharing wisdom (about Me) and always speaking of Me, My devotees derive contentment and delight.}
\slcol{\Index{तेषां सततयुक्तानां} भजतां प्रीतिपूर्वकम् । \\
ददामि बुद्धियोगं तं येन मामुपयान्ति ते ॥ १०.१० ॥ }
\translate{tēṣāṁ satatayuktānāṁ bhajatāṁ prītipūrvakam | \\
dadāmi buddhiyōgaṁ taṁ yēna māmupayānti tē || 10.10 || }
\cquote{To those who are ever steadfast, worship Me with love, I grant the \emph{yoga} of wisdom (\emph{buddhi-yogam}) which enables them to come to Me.}

\newpage
\begin{mananam}{\mananamfont {Mananam Śloka - 7}}
\normalfont\mananamtext How can learning about the various manifestations of divine power benefit me? Am I also endowed with this divine power such that I am able to manifest certain ideals in my life? How may I awaken, cultivate, and develop the Lord’s \emph{yogic} power within myself?
What unique talents do I have that I can start developing further? What impedes me from expressing my creativity through my work, hobbies, or other activities? Which areas of my life can I focus on to further develop my creative potential?
\end{mananam}
\sreflect
\begin{inspiration}{\mananamfont {Inspiration}}
\normalfont\mananamtext Just as the Divine possesses the ability to create and manifest a vast range of life expressions, humans too are endowed with the power of creativity. Those who learn to harness this ability are able to tap into their inherent divine potential. Sustaining positive thoughts and diligently working toward them is the foundation of success in any field. This ‘\emph{yoga}’—the practice of steadfastly holding onto one’s goals and ideals with divine faith and surrender, despite trials and challenges—serves as a bridge connecting us to the divine essence within ourselves.
\end{inspiration}
\newpage
\begin{mananam}{\mananamfont {Mananam Śloka - 8, 9}}
\normalfont\mananamtext What reasons and needs make seeking refuge in the Lord essential? In what ways can the Lord be the ultimate source of refuge and comfort?
Why might seeking refuge in power, money, and material goods be insufficient?
If we are finding support and comfort in family and friends, should we still prioritise seeking refuge in the Lord first? How can seeking refuge in the Lord—while encouraging our loved ones to do the same—be a way toward experiencing genuine fulfilment?
What distinguishes consciously and wisely seeking refuge in the Lord from doing so out of blind faith and ignorance of His true nature?
\end{mananam}
\sreflect
\begin{inspiration}{\mananamfont {Inspiration}}
\normalfont\mananamtext You cannot truly love someone you do not know well. Infatuation may seemingly mirror love, but those with maturity understand that infatuations are fleeting. Similarly, when we come to understand who the Lord truly is and how He is the very essence of our being, we naturally learn to turn to that One whose unwavering love and support remain the only constant in this ever-changing world. Those from whom we seek comfort—whether friends or family—may change their hearts or may not always be near us when we need them.
When our love for the Lord arises from genuine understanding, it brings a steady purpose to the mind. Love driven merely by impulse or fleeting emotion will inevitably fade as quickly as it appears.
\end{inspiration}
\newpage

\slcol{\Index{तेषामेवानुकम्पार्थम}हमज्ञानजं तमः । \\
नाशयाम्यात्मभावस्थो ज्ञानदीपेन भास्वता ॥ १०.११ ॥ }
\translate{tēṣāmēvānukampārthamahamajñānajaṁ tamaḥ | \\
nāśayāmyātmabhāvasthō jñānadīpēna bhāsvatā || 10.11 || }
\cquote{Out of mere compassion for them, I dwell in their hearts and destroy the darkness born of ignorance with the radiant lamp of knowledge.}
\slcol{अर्जुन उवाच —\\
\Index{परं ब्रह्म परं धाम} पवित्रं परमं भवान् । \\
पुरुषं शाश्वतं दिव्यमादिदेवमजं विभुम् ॥ १०.१२ ॥ }
\translate{arjuna uvāca —\\
paraṁ brahma paraṁ dhāma pavitraṁ paramaṁ bhavān | \\
puruṣaṁ śāśvataṁ divyamādidēvamajaṁ vibhum || 10.12 || }
\cquote{Arjuna said: You are the Supreme \emph{Brahman}, the Supreme Abode, the Supreme Purifier. You are the eternal Divine \emph{Puruṣa}, the primal God, unborn and omnipresent.}
\slcol{\Index{आहुस्त्वामृषयः सर्वे} देवर्षिर्नारदस्तथा । \\
असितो देवलो व्यासः स्वयं चैव ब्रवीषि मे ॥ १०.१३ ॥ }
\translate{āhustvāmr̥ṣayaḥ sarvē dēvarṣirnāradastathā | \\
asitō dēvalō vyāsaḥ svayaṁ caiva bravīṣi mē || 10.13 || }
\cquote{All the sages proclaim You as such; so too does the divine sage Nārada, as well as Asita, Devala, and Vyāsa. And now You Yourself declare it to me.}
\slcol{\Index{सर्वमेतदृतं मन्ये} यन्मां वदसि केशव । \\
न हि ते भगवन्व्यक्तिं विदुर्देवा न दानवाः ॥ १०.१४ ॥ }
\translate{sarvamētadr̥taṁ manyē yanmāṁ vadasi kēśava | \\
na hi tē bhagavanvyaktiṁ vidurdēvā na dānavāḥ || 10.14 || }
\cquote{I believe everything You say to me is true, O Keśava (Krishna). Neither the gods nor the demons comprehend Your true nature, O Lord.}

\newpage
\begin{mananam}{\mananamfont {Mananam Śloka - 10, 11}}
\normalfont\mananamtext Why do sages regard wisdom, in its purest and most unfiltered form, as the most sacred pursuit? Should we not, instead, seek special powers to manifest our aspirations and worthwhile dreams? If not, why not?  
What might be the outcome when our ignorance is completely destroyed? How does the wisdom of the Lord differ from the knowledge of nature’s workings, as part of God’s creation? What is my understanding of “\emph{buddhi-yoga}”?
\end{mananam}
\sreflect
\begin{inspiration}{\mananamfont {Inspiration}}
\normalfont\mananamtext The fruit of all spiritual practices and devotion to the Lord is the blessing of “\emph{buddhi-yoga}.” With this divine wisdom, we uproot the very cause of \emph{saṁsāra}— the intricate web of worldly existence that has kept us bound for countless lifetimes. The darkness of ignorance is so profound that the Lord, who is the very essence of our being and dwells within us, remains unnoticed. When our spiritual practices reach their peak, our devotion becomes pure, and our discernment of reality—honed by the guidance of the \emph{guru} and the \emph{śāstras}—grows razor-sharp. It is then that we receive the grace of the Lord, culminating in God-realisation.
\end{inspiration}
\newpage

\slcol{\Index{स्वयमेवात्मनात्मानं} वेत्थ त्वं पुरुषोत्तम । \\
भूतभावन भूतेश देवदेव जगत्पते ॥ १०.१५ ॥ }
\translate{svayamēvātmanātmānaṁ vēttha tvaṁ puruṣōttama | \\
bhūtabhāvana bhūtēśa dēvadēva jagatpatē || 10.15 || }
\cquote{Indeed, You alone know Your true nature through Your own Self-awareness, O \emph{Puruṣottama} (Supreme Person), \emph{Bhūtabhāvana} (Origin of all beings), \emph{Bhūteśa} (Lord of all beings), \emph{Devadeva} (God of gods) and \emph{Jagatpate} (Lord of the universe).}
\slcol{\Index{वक्तुमर्हस्यशेषेण} दिव्या ह्यात्मविभूतयः । \\
याभिर्विभूतिभिर्लोकानिमांस्त्वं व्याप्य तिष्ठसि ॥ १०.१६ ॥ }
\translate{vaktumarhasyaśēṣēṇa divyā hyātmavibhūtayaḥ | \\
yābhirvibhūtibhirlōkānimāṁstvaṁ vyāpya tiṣṭhasi || 10.16 || }
\cquote{Please describe to me in detail Your divine glories, by which You pervade all these worlds and reside in them.}
\slcol{\Index{कथं विद्यामहं योगिंस्त्वां} सदा परिचिन्तयन् । \\
केषु केषु च भावेषु चिन्त्योऽसि भगवन्मया ॥ १०.१७ ॥ }
\translate{kathaṁ vidyāmahaṁ yōgiṁstvāṁ sadā paricintayan | \\
kēṣu kēṣu ca bhāvēṣu cintyō:'si bhagavanmayā || 10.17 || }
\cquote{O \emph{Yogin} (Krishna) how may I, ever-meditating, know You? In what various forms (or aspects) are You to be contemplated by me, O Blessed Lord?}
\slcol{\Index{विस्तरेणात्मनो योगं} विभूतिं च जनार्दन । \\
भूयः कथय तृप्तिर्हि शृण्वतो नास्ति मेऽमृतम् ॥ १०.१८ ॥ }
\translate{vistarēṇātmanō yōgaṁ vibhūtiṁ ca janārdana | \\
bhūyaḥ kathaya tr̥ptirhi śr̥ṇvatō nāsti mē:'mr̥tam || 10.18 || }
\cquote{O Janārdana (maintainer of all living beings, Krishna), please describe again in detail Your divine glories and \emph{yoga}; for I am never satisfied while hearing Your nectar-like words.}

\newpage
\begin{mananam}{\mananamfont {Mananam Śloka - 17, 18}}
\normalfont\mananamtext Is the Supreme Divine Being confined to a specific form, like humans? What would happen if God were limited in power and capacity? 
Could it be beneficial to understand God beyond a single form? How can I discover which image or aspect of the Divine is most suitable for me to meditate upon? 
Would it be helpful to meditate upon or pray to different aspects of God based on my changing mental states? Can I use my God-given creativity to visualise and conceptualise a divine image that resonates with me personally?
\end{mananam}
\sreflect
\begin{inspiration}{\mananamfont {Inspiration}}
\normalfont\mananamtext Our human understanding of God will always be limited, as only the Lord can fully reveal His true nature. While it is helpful to seek the Lord in a personal, human form, He is not confined to any single form. The grandest, purest, and most sublime aspects of all creation reflect the presence of the Lord. Therefore, in every moment of conscious living, we can find opportunities to remember and contemplate Him.
\end{inspiration}
\newpage

\slcol{श्रीभगवानुवाच —\\
\Index{हन्त ते कथयिष्यामि} दिव्या ह्यात्मविभूतयः । \\
प्राधान्यतः कुरुश्रेष्ठ नास्त्यन्तो विस्तरस्य मे ॥ १०.१९ ॥ }
\translate{śrībhagavānuvāca —\\
hanta tē kathayiṣyāmi divyā hyātmavibhūtayaḥ | \\
prādhānyataḥ kuruśrēṣṭha nāstyantō vistarasya mē || 10.19 || }
\cquote{Sri Bhagavan said: Indeed, I shall now describe to you My divine glories, O \emph{kuru-śhreṣhṭha} (best of the Kurus, Arjuna), but only the most prominent ones—for there is no end to the extent of My manifestations.}
\slcol{\Index{अहमात्मा गुडाकेश} सर्वभूताशयस्थितः । \\
अहमादिश्च मध्यं च भूतानामन्त एव च ॥ १०.२० ॥ }
\translate{ahamātmā guḍākēśa sarvabhūtāśayasthitaḥ | \\
ahamādiśca madhyaṁ ca bhūtānāmanta ēva ca || 10.20 || }
\cquote{O Guḍākeśha (conqueror of sleep, Arjuna), I am the Self dwelling in the hearts of all beings. I am the beginning, the middle, and the end of all beings.}
\slcol{\Index{आदित्यानामहं विष्णुर्ज्योतिषां} रविरंशुमान् । \\
मरीचिर्मरुतामस्मि नक्षत्राणामहं शशी ॥ १०.२१ ॥ }
\translate{ādityānāmahaṁ viṣṇurjyōtiṣāṁ raviraṁśumān | \\
marīcirmarutāmasmi nakṣatrāṇāmahaṁ śaśī || 10.21 || }
\cquote{Among the \emph{Ādityas} (twelve deities representing different aspects of the sun) I am Viṣhṇu; among luminaries I am the radiant sun. Among the \emph{Maruts} (storm-gods) I am Marīchi, and among the stars I am the moon.}
\slcol{\Index{वेदानां सामवेदोऽस्मि} देवानामस्मि वासवः । \\
इन्द्रियाणां मनश्चास्मि भूतानामस्मि चेतना ॥ १०.२२ ॥ }
\translate{vēdānāṁ sāmavēdō:'smi dēvānāmasmi vāsavaḥ | \\
indriyāṇāṁ manaścāsmi bhūtānāmasmi cētanā || 10.22 || }
\cquote{Among the vedas I am the \emph{Sāma veda}; of the gods I am \emph{Vāsava} (Indra, king of heaven). Among the senses I am the mind, and among living beings I am \emph{cetanā}, the sentient principle.}

\newpage

\begin{mananam}{\mananamfont {Mananam Śloka - 20}}
\normalfont\mananamtext How can I perceive God dwelling within my heart? Does this refer to the physical heart or the spiritual heart? In what way can God be my very Self? What is the nature of my innermost Self?
When it is said that God is the beginning of all beings, does it refer to the act of world creation? If God is also present in the middle and end of all beings, how can I comprehend such an uninterrupted stream of existence?
\end{mananam}
\sreflect
\begin{inspiration}{\mananamfont {Inspiration}}
\normalfont\mananamtext The very essence of our being, the pure awareness we identify as \enquote{I am} without any further labels, is the Supreme Lord Himself. This supreme, pure existence is what enables us to become this conscious individual and experience life. The core of all knowing and beingness is the Lord alone. As finite beings, we arise from the vast ocean of the Lord’s existence, express ourselves as individual waves of life, and ultimately merge back into that infinite ocean.
\end{inspiration}

\newpage
\begin{mananam}{\mananamfont {Mananam Śloka - 21, 22}}
\normalfont\mananamtext Does the sun’s brilliance not surpass all other sources of light? Does the moon not outshine all the stars? Meditate on God as the brightest and most radiant among all shining entities.
Can my senses operate without the mind? Can my mind function without intelligence? Can the various aspects of my body and mind work without consciousness? Reflect on that higher intelligence—the unifying force behind the functioning of the senses, mind, and intellect.
\end{mananam}
\sreflect
\begin{inspiration}{\mananamfont {Inspiration}}
\normalfont\mananamtext Among all natural lights—fire, the moon, stars, and lightning—the sun shines the brightest. God can be meditated upon as the supreme light in this world, yet ultimately, He is the Light of all lights.\\
In the same way, God is the light of consciousness. It is His divine light that illuminates the mind, enabling it to receive input from the senses and empowering the intelligence to make decisions.
\end{inspiration}
\newpage

\slcol{\Index{रुद्राणां शङ्करश्चास्मि} वित्तेशो यक्षरक्षसाम् । \\
वसूनां पावकश्चास्मि मेरुः शिखरिणामहम् ॥ १०.२३ ॥ }
\translate{rudrāṇāṁ śaṅkaraścāsmi vittēśō yakṣarakṣasām | \\
vasūnāṁ pāvakaścāsmi mēruḥ śikhariṇāmaham || 10.23 || }
\cquote{Among the \emph{Rudras}, I am Śaṅkara (Śiva); among the \emph{Yakṣas} and \emph{Rākṣasas}, I am Kubera, the lord of wealth; among the \emph{Vasus}, I am Agni (fire); and among mountains, I am Meru, the highest peak.}
\slcol{\Index{पुरोधसां च मुख्यं मां} विद्धि पार्थ बृहस्पतिम् । \\
सेनानीनामहं स्कन्दः सरसामस्मि सागरः ॥ १०.२४ ॥ }
\translate{purōdhasāṁ ca mukhyaṁ māṁ viddhi pārtha br̥haspatim | \\
sēnānīnāmahaṁ skandaḥ sarasāmasmi sāgaraḥ || 10.24 || }
\cquote{Among priests, O Pārtha, know Me to be the chief, Bṛihaspati (the celestial priest). Among army commanders I am Skanda (the god of war, Karthikeya), and among reservoirs of water I am the ocean.}
\slcol{\Index{महर्षीणां भृगुरहं} गिरामस्म्येकमक्षरम् । \\
यज्ञानां जपयज्ञोऽस्मि स्थावराणां हिमालयः ॥ १०.२५ ॥ }
\translate{maharṣīṇāṁ bhr̥gurahaṁ girāmasmyēkamakṣaram | \\
yajñānāṁ japayajñō:'smi sthāvarāṇāṁ himālayaḥ || 10.25 || }
\cquote{Among the great sages I am Bhṛigu; Among words I am the syllable \emph{Om}. Among sacrifices I am \emph{japa} (silent repetition of the holy names), and among immovable things I am the Himalayas.}
\slcol{\Index{अश्वत्थः सर्ववृक्षाणां} देवर्षीणां च नारदः । \\
गन्धर्वाणां चित्ररथः सिद्धानां कपिलो मुनिः ॥ १०.२६ ॥ }
\translate{aśvatthaḥ sarvavr̥kṣāṇāṁ dēvarṣīṇāṁ ca nāradaḥ | \\
gandharvāṇāṁ citrarathaḥ siddhānāṁ kapilō muniḥ || 10.26 || }
\cquote{Among all trees, I am the \emph{Aśvattha} (sacred fig); among the Devarṣis, I am Nārada; among the Gandharvas, I am Citraratha; and among the Siddhas, I am the sage Kapila.}

\newpage
\begin{mananam}{\mananamfont {Mananam Śloka - 25}}
\normalfont\mananamtext Among the three \emph{guṇas} in nature, which can be considered the most fitting representation of divinity in holy ones and sages?
Why is \emph{Oṃ} regarded as the most sacred of all syllables? Is it purely a matter of faith and tradition, or could there be a deeper mental and spiritual significance to it?
How can \emph{japa} (the repetition of a \emph{mantra}) be considered a form of \emph{yajña} (sacrifice)? What exactly are we offering when silently repeating a \emph{mantra}?
Why did the Lord choose the Himalayas to symbolize stillness? How can the physical stillness of a mountain influence my mind? 
\end{mananam}
\sreflect
\begin{inspiration}{\mananamfont {Inspiration}}
\normalfont\mananamtext In this \emph{śloka}, the Supreme Being is meditated upon through the manifestations of \emph{Ṛṣi} Bhrigu, \emph{japa}, the sacred syllable \emph{Om}, and the Himalayas.
If all the \emph{ṛṣis} embody \emph{sattvic} energy, then God is the purest and highest expression of \emph{sattva}. Likewise, the ancient practice of \emph{japa}, the chanting of \emph{Om}, and the stillness of the Himalayas each serve as pathways to connect with the divine presence within us.
\end{inspiration}
\newpage

\slcol{\Index{उच्चैःश्रवसमश्वानां} विद्धि माममृतोद्भवम् । \\
ऐरावतं गजेन्द्राणां नराणां च नराधिपम् ॥ १०.२७ ॥ }
\translate{uccaiḥśravasamaśvānāṁ viddhi māmamr̥tōdbhavam | \\
airāvataṁ gajēndrāṇāṁ narāṇāṁ ca narādhipam || 10.27 || }
\cquote{Among horses, know Me as Uccaiḥśravas, born of \emph{amṛta} (nectar); among elephants, I am Airāvata; and among men, I am the king.}
\slcol{\Index{आयुधानामहं वज्रं} धेनूनामस्मि कामधुक् । \\
प्रजनश्चास्मि कन्दर्पः सर्पाणामस्मि वासुकिः ॥ १०.२८ ॥ }
\translate{āyudhānāmahaṁ vajraṁ dhēnūnāmasmi kāmadhuk | \\
prajanaścāsmi kandarpaḥ sarpāṇāmasmi vāsukiḥ || 10.28 || }
\cquote{Among weapons, I am the Vajra (thunderbolt); among cows, I am Kāmadhuk (the wish-fulfilling cow); among causes of procreation, I am Kandarpa (the god of love); and among serpents, I am Vāsuki.}
\slcol{\Index{अनन्तश्चास्मि नागानां} वरुणो यादसामहम् । \\
पितॄणामर्यमा चास्मि यमः संयमतामहम् ॥ १०.२९ ॥ }
\translate{anantaścāsmi nāgānāṁ varuṇō yādasāmaham | \\
pitr̥̄ṇāmaryamā cāsmi yamaḥ saṁyamatāmaham || 10.29 || }
\cquote{Among the \emph{Nāgas}, I am Ananta; among the \emph{Yādasas} (aquatic beings), I am Varuṇa; among the \emph{pitṛs} (ancestors), I am Aryamā; and among the \emph{saṁyamatām} (controllers), I am Yama.}
\slcol{\Index{प्रह्लादश्चास्मि दैत्यानां} कालः कलयतामहम् । \\
मृगाणां च मृगेन्द्रोऽहं वैनतेयश्च पक्षिणाम् ॥ १०.३० ॥ }
\translate{prahlādaścāsmi daityānāṁ kālaḥ kalayatāmaham | \\
mr̥gāṇāṁ ca mr̥gēndrō:'haṁ vainatēyaśca pakṣiṇām || 10.30 || }
\cquote{Among the \emph{Daityas}, I am \emph{Prahlāda}; among the \emph{Kalayatām} (reckoners of time), I am \emph{Kāla} (time); among the \emph{Mṛgāṇām} (beasts), I am the \emph{Mṛgendra} (lion); and among the \emph{Pakṣiṇām} (birds), I am \emph{Vainateya} (Garuḍa).}

\newpage
\slcol{\Index{पवनः पवतामस्मि} रामः शस्त्रभृतामहम् । \\
झषाणां मकरश्चास्मि स्रोतसामस्मि जाह्नवी ॥ १०.३१ ॥ }
\translate{pavanaḥ pavatāmasmi rāmaḥ śastrabhr̥tāmaham | \\
jhaṣāṇāṁ makaraścāsmi srōtasāmasmi jāhnavī || 10.31 || }
\cquote{Among purifiers I am the Pavana (Wind); among wielders of weapons, I am Lord Rāma. Among aquatic creatures I am the Shark, and among flowing rivers I am Jāhnavī (daughter of Jahnu, Ganges).}
\slcol{\Index{सर्गाणामादिरन्तश्च} मध्यं चैवाहमर्जुन । \\
अध्यात्मविद्या विद्यानां वादः प्रवदतामहम् ॥ १०.३२ ॥ }
\translate{sargāṇāmādirantaśca madhyaṁ caivāhamarjuna | \\
adhyātmavidyā vidyānāṁ vādaḥ pravadatāmaham || 10.32 || }
\cquote{Of the \emph{sargāṇām} (creations), I am the beginning, the end, and also the middle, O Arjuna; among knowledges, I am the \emph{adhyātmavidyā}, Self-knowledge; and among the \emph{pravadatām} (disputants), I am \emph{vāda} (reasoned debate).}
\slcol{\Index{अक्षराणामकारोऽस्मि} द्वन्द्वः सामासिकस्य च । \\
अहमेवाक्षयः कालो धाताहं विश्वतोमुखः ॥ १०.३३ ॥ }
\translate{akṣarāṇāmakārō:'smi dvandvaḥ sāmāsikasya ca | \\
ahamēvākṣayaḥ kālō dhātāhaṁ viśvatōmukhaḥ || 10.33 || }
\cquote{Among the alphabets, I am \emph{akāra}, the letter ‘A’; among the compounds, I am the \emph{dvandva}, dual compound; I alone am \emph{kāla}, eternal time; and I am the \emph{dhātā}, dispenser (of fruits of actions), having faces in all directions. }
\slcol{\Index{मृत्युः सर्वहरश्चाहमुद्भवश्च} भविष्यताम् । \\
कीर्तिः श्रीर्वाक्च नारीणां स्मृतिर्मेधा धृतिः क्षमा ॥ १०.३४ ॥ }
\translate{mr̥tyuḥ sarvaharaścāhamudbhavaśca bhaviṣyatām | \\
kīrtiḥ śrīrvākca nārīṇāṁ smr̥tirmēdhā dhr̥tiḥ kṣamā || 10.34 || }
\cquote{I am all-devouring death, and I am the origin of those yet to be born (alternatively—[I am] the prosperity of those who are to be prosperous).  Among feminine qualities, I am fame, prosperity, fine speech, memory, intelligence, steadfastness, and forbearance. }

\newpage
\begin{mananam}{\mananamfont {Mananam Śloka - 31, 32}}
\normalfont\mananamtext Why is wind regarded as a purifying agent? How can this cleansing quality of wind be seen as a representation of the Lord for meditation? 
What makes the Ganga so sacred that it is chosen as the ideal representation of the Divine among all rivers? 
Why is Self-knowledge considered the most precious form of wisdom? How does the knowledge of the Self relate to the knowledge of the Lord? Are these two separate truths to be known and realised, or are they interconnected?
\end{mananam}
\sreflect
\begin{inspiration}{\mananamfont {Inspiration}}
\normalfont\mananamtext The passing wind and the flowing river are both manifestations of the Lord in dynamic activity. To be active means constant movement and change. If nature were stagnant, it would lack the richness and vitality of creation. Yet, recognising the unchanging Lord within the ever-changing flow of nature brings deep solace to the devotee’s heart.
Seeking the Self is the highest pursuit of human life. Of what value is the knowledge of countless things if one remains unaware of the very source that makes all “knowing” and “learning” possible? To realise “Who am I?” is not only the pinnacle of wisdom but the only knowledge that quiets restless curiosity—the force that drives us outward, leaving us perpetually unsatisfied.
\end{inspiration}

\newpage
\begin{mananam}{\mananamfont {Mananam Śloka - 34}}
\normalfont\mananamtext Why is death described as that which takes everything away? How can this phenomenon represent a manifestation of the all-powerful divine presence in creation?
How can the same Lord who takes everything away also be the one who bestows prosperity? What is the significance of potentiality, expressed through the use of the future tense, in the context of bestowing prosperity? 
Why are qualities such as fame, wealth, memory, intelligence, firmness, and forgiveness considered feminine? Does this imply that men cannot aspire to cultivate these virtues?
\end{mananam}
\sreflect
\begin{inspiration}{\mananamfont {Inspiration}}
\normalfont\mananamtext Nearly all humans fear authorities and laws that can exercise power over us. While such power, whether legal or illegal, can strip us of wealth, possessions, and freedom, the highest power is the one capable of taking life itself.
Likewise, the ultimate power is also the pinnacle of opulence. The Lord, as the giver of life, holds the potential to bless even those who are poor and deprived of worldly possessions with prosperity and abundance.
These \emph{mantras} serve as a profound meditation on the glory of the Lord—not only as the one who can take life away but also as the giver who breathes new life into a mother’s womb or causes a seed to sprout into a plant or tree.
\end{inspiration}
\newpage

\slcol{\Index{बृहत्साम तथा साम्नां} गायत्री च्छन्दसामहम् । \\
मासानां मार्गशीर्षोऽहमृतूनां कुसुमाकरः ॥ १०.३५ ॥ }
\translate{br̥hatsāma tathā sāmnāṁ gāyatrī cchandasāmaham | \\
māsānāṁ mārgaśīrṣō:'hamr̥tūnāṁ kusumākaraḥ || 10.35 || }
\cquote{Among the hymns in the \emph{Sāma veda}, I am the \emph{bṛihat-sāma} (a specific hymn in the \emph{Sāma veda}); of all poetic meters, I am the \emph{gāyatrī}; of months, I am \emph{mārgaśhīrṣha} (approx. November–December); of seasons, I am spring, the season of flowers.}
\slcol{\Index{द्यूतं छलयतामस्मि} तेजस्तेजस्विनामहम् । \\
जयोऽस्मि व्यवसायोऽस्मि सत्त्वं सत्त्ववतामहम् ॥ १०.३६ ॥ }
\translate{dyūtaṁ chalayatāmasmi tējastējasvināmaham | \\
jayō:'smi vyavasāyō:'smi sattvaṁ sattvavatāmaham || 10.36 || }
\cquote{I am the gambling of the deceitful, and the splendour of the splendid. I am victory, I am determination, and I am the \emph{sattva} quality of those possessed of \emph{sattva}.}
\slcol{\Index{वृष्णीनां वासुदेवोऽस्मि} पाण्डवानां धनञ्जयः । \\
मुनीनामप्यहं व्यासः कवीनामुशना कविः ॥ १०.३७ ॥ }
\translate{vr̥ṣṇīnāṁ vāsudēvō:'smi pāṇḍavānāṁ dhanañjayaḥ | \\
munīnāmapyahaṁ vyāsaḥ kavīnāmuśanā kaviḥ || 10.37 || }
\cquote{Among the descendants of \emph{Vṛṣhṇī}, I am Vāsudeva (son of Vasudeva, Krishna); among the \emph{Pāṇḍavās}, I am Dhanañjaya (the winner of wealth, Arjuna). Among the sages, I am Vyāsa, and among great thinkers, I am Uśhanā.}
\slcol{\Index{दण्डो दमयतामस्मि} नीतिरस्मि जिगीषताम् । \\
मौनं चैवास्मि गुह्यानां ज्ञानं ज्ञानवतामहम् ॥ १०.३८ ॥ }
\translate{daṇḍō damayatāmasmi nītirasmi jigīṣatām | \\
maunaṁ caivāsmi guhyānāṁ \\
\hspace*{0.5cm}jñānaṁ jñānavatāmaham || 10.38 || }
\cquote{Among those who enforce discipline, I am the rod of punishment (instrument of discipline); among those who seek victory, I am statesmanship (morality); of secrets, I am silence (\emph{mauna}), and of the wise, I am wisdom (\emph{jñāna}) itself.}

\newpage
\begin{mananam}{\mananamfont {Mananam Śloka - 36}}
\normalfont\mananamtext How can the Supreme Being manifest as manipulation and deceit among gamblers? Isn’t it contradictory that the one worshipped and adored in a personal form can also express such undesirable qualities?
Can the God we revere, while displaying such qualities, remain pure and untainted by them?
How is the Lord the manifestation of victory? Isn’t victory often associated with glory and achievement?
How is the Lord the essence of goodness and purity?
What in my worldly experiences can I identify as the highest expression of goodness? Can I use that experience or entity as a guide for meditating on the Lord?
\end{mananam}
\sreflect
\begin{inspiration}{\mananamfont {Inspiration}}
\normalfont\mananamtext When the Supreme Spirit manifests in nature, it can take on any extreme form (for all forms and qualities ultimately belong to the manifest God). Only the Supreme can be powerful enough to counter and overcome even the most deceitful and destructive of forces. The scriptures teach that among all illusions, \emph{maya} plays the greatest trick, for it veils the reality of the Lord from our awareness.
Though the pure Spirit can manifest as both good and bad, it remains valuable for a seeker to revere the Lord as the pinnacle of goodness and purity. By doing so, one cultivates and expresses those virtues within oneself. \emph{Sattva} (purity) leads toward freedom and liberation, while \emph{tamas} (inertia) binds us more deeply to ignorance. Both arise from the Lord’s nature but do not stain or diminish the essence of the Divine.
\end{inspiration}
\newpage

\begin{mananam}{\mananamfont {Mananam Śloka - 38}}
\normalfont\mananamtext Is it possible that God administers justice and punishes those who harm the innocent and pious? Do I believe that the divine \emph{avatāras} who manifested in order to destroy evil and wickedness were real historical figures, or do I see them more as symbolic stories?
Why is silence considered the greatest secret? Can it reveal truths that are too profound or ineffable to be expressed in words?
What is the highest form of knowledge that one can possess? Can a person study and retain all knowledge of the material world throughout life?
Is there an essential wisdom that serves as the key to understanding all fields of study?
\end{mananam}
\sreflect
\begin{inspiration}{\mananamfont {Inspiration}}
\normalfont\mananamtext Nearly all religions uphold a belief in some form of divine retribution. Whether through a direct act of God or the precise workings of \emph{karmic} law, those who persist in evil intentions and actions inevitably face the consequences of their deeds. Therefore, it is wiser to seek justice through God or the divine force governing creation.
\emph{Munis} are those who have mastered control over their speech and minds. In the profound silence they cultivate, they uncover the deeper mysteries of creation and the realms beyond.
True knowledge encompasses the understanding of God, creation, and, most importantly, our own being. It is this essential knowingness—the subjective awareness within—that empowers us to perceive and comprehend all objective knowledge.
\end{inspiration}
\newpage

\slcol{\Index{यच्चापि सर्वभूतानां} बीजं तदहमर्जुन । \\
न तदस्ति विना यत्स्यान्मया भूतं चराचरम् ॥ १०.३९ ॥ }
\translate{yaccāpi sarvabhūtānāṁ bījaṁ tadahamarjuna | \\
na tadasti vinā yatsyānmayā bhūtaṁ carācaram || 10.39 || }
\cquote{Furthermore, O Arjuna, I am the \emph{bījaṁ}, seed, of all beings. There is no being—whether moving or non-moving—that can exist without Me.}
\slcol{\Index{नान्तोऽस्ति मम दिव्यानां} विभूतीनां परन्तप । \\
एष तूद्देशतः प्रोक्तो विभूतेर्विस्तरो मया ॥ १०.४० ॥ }
\translate{nāntō:'sti mama divyānāṁ vibhūtīnāṁ parantapa | \\
ēṣa tūddēśataḥ prōktō vibhūtērvistarō mayā || 10.40 || }
\cquote{O Parantapa (conqueror of enemies, Arjuna) there is no end to My \emph{divyānām vibhūtīnām}, divine manifestations. What I have described to you is but a brief account of My \emph{vibhūti}, glories.}
\slcol{\Index{यद्यद्विभूतिमत्सत्त्वं} श्रीमदूर्जितमेव वा । \\
तत्तदेवावगच्छ त्वं मम तेजोंशसम्भवम् ॥ १०.४१ ॥ }
\translate{yadyadvibhūtimatsattvaṁ śrīmadūrjitamēva vā | \\
tattadēvāvagaccha tvaṁ mama tējōṁśasambhavam || 10.41 || }
\cquote{Know that whatever is glorious, prosperous, or mighty in this world is but a portion of My own splendour.}
\slcol{\Index{अथवा बहुनैतेन} किं ज्ञातेन तवार्जुन । \\
विष्टभ्याहमिदं कृत्स्नमेकांशेन स्थितो जगत् ॥ १०.४२ ॥ }
\translate{athavā bahunaitēna kiṁ jñātēna tavārjuna | \\
viṣṭabhyāhamidaṁ kr̥tsnamēkāṁśēna sthitō jagat || 10.42 || }
\cquote{But what is the need for all this detailed knowledge, O Arjuna? Simply know that I pervade and support this entire universe with just one part of Myself.}

\newpage
\begin{mananam}{\mananamfont {Mananam Śloka - 39, 40}}
\normalfont\mananamtext If God is the seed of all beings, does that mean that everything in creation is a manifestation of the Divine? Can I see the presence of God in every aspect of nature and life around me?
Doesn’t all matter have a traceable origin? Just as science identifies atoms as the fundamental units of matter, could there also be a divine essence as the ultimate source of everything and everyone?
Has human science been able to trace the origin of the energy that keeps matter in motion or sustains it in stillness? Is it possible to keep tracing the source of creation in a hierarchical order? Will there ever be a definitive end to this search for the ultimate source?
\end{mananam}
\sreflect
\begin{inspiration}{\mananamfont {Inspiration}}
\normalfont\mananamtext Even what science describes as energy or the initial spark of the Big Bang must have its origin in some preceding God-particle or primal source. The human mind, with its inherent limitations, tends to seek finite units to define the infinite. Yet, the true essence of God transcends all such boundaries and cannot be confined to intellectual definitions.
The glories and manifestations of the Divine are beyond the grasp of the ordinary human mind. Only when the intellect is elevated to a divine state—when it merges with higher wisdom—can the profound mysteries of matter, energy, and existence be fully revealed.
\end{inspiration}
\newpage
\begin{mananam}{\mananamfont {Mananam Śloka - 41, 42}}
\normalfont\mananamtext Whenever I see someone powerful, talented, or rich, do I realise that it is a manifestation of the Divine potential? Time has shown that despite one human who may be great and special, someone else comes along and shows a higher greatness. Have I considered why world records get broken frequently? Is there a record that will never be broken? 
Is it possible that God is the epitome of all that we admire in people and in nature? Is it also possible that all the talent, power and riches in the world are not anyone’s own, but borrowed from the Supreme Being? Can this Supreme Being be the essential basis for everything in this universe?
\end{mananam}
\sreflect
\begin{inspiration}{\mananamfont {Inspiration}}
\normalfont\mananamtext All the great qualities and talents we admire in people, all the powers that one might have, all the riches that one may possess, are all a manifestation of the Supreme Being. In the Lord’s creation, there are unlimited and undreamed-of possibilities. All of the universes that we know speculate about are only a small portion of the Divine Being. Without God, nothing exists. A wise devotee is one who learns to use every spectacle as a meditation pointer toward the Lord. Furthermore, the devotee realises that deep within is that supremely creative and manifesting force of the Divine to transform all the unhealthy and negative conditions in life. 
\end{inspiration}

\newpage
\begin{center}
\devfont{इति श्रीमद्भगवद्नीतासूपनिषत्सु  ब्रह्मविद्यायां  योगशास्त्रे
श्रीकृष्णार्जुनसंवादे विभूतियोगो  नाम दशमोऽध्यायः  ॥ }\\
\chapEndSlokaTranslate{ōṃ tatsaditi śrīmadbhagavadgītāsūpaniṣatsu brahmavidyāyāṃ 
yōgaśāstrē śrīkṛṣṇārjunasaṃvādē vibhūtiyōgō nāma daśamō'dhyāyaḥ ||}
\end{center}